\originalchapter{习题课七: 下落线圈}
题源为\href{https://ocw.mit.edu/courses/8-02-physics-ii-electricity-and-magnetism-spring-2007/resources/solving07/}{Problem Solving 7}, PDF第1--6页(原页脚沿用Solving 8). 全部Q1--Q10保留, 所有解答为AI独立解答.

\begin{exercise}
质量$m$、宽$w$、竖直高$\ell$、电阻$R$的矩形线圈在重力下向下落出磁场区. 场区内$\bb=B\uv x$指向纸外, 场区外为零. 当下边已在场外、上边仍在场内时, 上边距场区下边界$z(t)>0$, 正$z$竖直向上, 速度$V(t)=\dot z<0$. 选回路面积法向为纸外. 忽略线圈自感和边界过渡宽度.
Q1求$V$与$z$的关系.
Q2用$z$写磁通.
Q3用$V$写通量变化率并判正负.
Q4判电流顺逆时针.
Q5判感生场方向, 是否阻碍原磁通改变?
Q6求电流大小.
Q7求除重力外竖直磁力大小和方向.
Q8求终端速率大小.
Q9写终端时重力功率与焦耳热功率相等的通用关系.
Q10代入核对.
\end{exercise}
\begin{solution}
Q1. $V=\dot z<0$. Q2. 场区内面积$wz$, 所以$\Phi_B=Bwz$.
Q3. $\dot\Phi_B=BwV<0$.
Q4. $\mathcal E=-BwV>0$, 相对于纸外法向右手环向即逆时针.
Q5. 逆时针电流在环内产生纸外场, 补偿纸外通量的减少, 不是使磁场方向反转.
Q6. $I=-BwV/R=Bw|V|/R$.
Q7. 两竖边的水平力抵消; 下边在场外, 上边在场内电流向左. 左向电流与纸外场叉积向上, 得$F_z=IBw=-B^2w^2V/R>0$. 整体运动式为$m\dot V=-mg-(B^2w^2/R)V$.
Q8. 若在仍部分离开场的时间内达到准终端态, $V_t=-mgR/(B^2w^2)$, 速率$|V_t|=mgR/(B^2w^2)$. 线圈完全离开后磁力消失, 不会永远保持这一终端速度; 有限场区和有限线圈也未必来得及达到它.
Q9. $\vct F_g\cdot\vct V=I^2R$, 即$mg|V_t|=I_t^2R$.
Q10. 左边为$m^2g^2R/(B^2w^2)$. 右边用$I_t=Bw|V_t|/R=mg/(Bw)$也得相同值. 一般暂态中$mg|V|=I^2R+\dd(mV^2/2)/\dd t$, 终端时动能不变才可删去后一项.
\end{solution}
